% ======================================================================
%  APÉNDICE A — EL PROGRAMA DEL JUGUETE
% ======================================================================

\appendix
\chapter{El programa del juguete: \texttt{toy.py}}\label{app:codigo}

\begin{roadmap}
Este apéndice reproduce, íntegro y ejecutable, el programa que generó
las cifras de los laboratorios de los
capítulos~\ref{ch:labgoppa} y~\ref{ch:labkem}: el McEliece de juguete
sobre $\FF4$ con $n=16$ y $t=2$. El listado no es decorativo: cada
número impreso en esos capítulos está aquí calculado, y cada
afirmación delicada está protegida por una aserción que detiene el
programa si algo no cuadra.
\end{roadmap}

\section{Qué es y cómo se ejecuta}

\texttt{toy.py} es un único archivo de Python~3, sin más dependencia
que la biblioteca estándar (\texttt{hashlib.shake\_256}); ocupa unas
$450$ líneas. Se ejecuta con \texttt{python3 toy.py} y termina
imprimiendo \texttt{TODAS LAS COMPROBACIONES OK} si todo cuadra; si
una aserción falla, se detiene en la primera discrepancia. Reproduce,
en orden:

\begin{enumerate}[leftmargin=1.8em,itemsep=2pt]
  \item el cuerpo $\FF4$ y las $16$ columnas de la matriz $\tilde H$
        de la tabla~\ref{tab:lab-columnas} (capítulo~\ref{ch:labgoppa});
  \item la forma sistemática $(\mathrm{I}_8\,|\,T)$, la matriz
        $G=(T^{\top}\,|\,\mathrm{I}_8)$ y la permutación
        $\sigma$ del laboratorio (figura~\ref{fig:lab-SPG});
  \item el síndrome del error $e$ con soporte $\{10,13\}$ y la
        traza completa de Patterson, incluidos sus certificados;
  \item la comprobación de tipo van Lint--MacWilliams: las $2^{16}$
        palabras del código, contadas dos veces y comparadas;
  \item los $120$ errores de peso $\le 2$, los $137$ síndromes
        alcanzables y la lista completa de síndromes;
  \item la cadena del capítulo~\ref{ch:labkem}: la semilla $s$, el
        error $e$, el criptograma $C$, la clave $K$ y el criptograma
        manipulado $\tilde C$ con su rechazo implícito $K''$.
\end{enumerate}

\section{El listado completo}

\begin{lstlisting}
#!/usr/bin/env python3
# ----------------------------------------------------------------------
# toy.py - el juguete de bolsillo de los capitulos 4, 5 y 11
# Reproduce, paso a paso, todas las cifras impresas del libro.
# Uso:  python3 toy.py
# ----------------------------------------------------------------------

from hashlib import shake_256

# ======================================================================
# 1. GF(16) = F2[z]/(z^4+z+1)
# Cada elemento es un entero 0..15 (un digito hexadecimal): el bit l
# es el coeficiente de z^l.  z = 2, y z^4+1 = z.
# ======================================================================

def _tablas():
    exp, log = [0] * 15, [0] * 16
    v = 1
    for e in range(15):
        exp[e] = v
        log[v] = e
        v <<= 1
        if v & 0x10:
            v ^= 0x13                      # z^4 = z+1
    return exp, log

EXP, LOG = _tablas()

def mul(a, b):
    return 0 if not a or not b else EXP[(LOG[a] + LOG[b]) % 15]

def inv(a):
    return EXP[(-LOG[a]) % 15]

def pw(a, k):
    return 0 if not a else EXP[(LOG[a] * k) % 15]

# ======================================================================
# 2. Polinomios sobre GF(16): listas de coeficientes, indice = grado
# ======================================================================

def trim(p):
    p = p[:]
    while len(p) > 1 and p[-1] == 0:
        p.pop()
    return p

def p_add(a, b):
    r = [0] * max(len(a), len(b))
    for i, c in enumerate(a):
        r[i] ^= c
    for i, c in enumerate(b):
        r[i] ^= c
    return trim(r)

def p_esc(c, p):                           # multiplica por el escalar c
    return trim([mul(c, x) for x in p])

def p_mul(a, b):
    r = [0] * (len(a) + len(b) - 1)
    for i, x in enumerate(a):
        if x:
            for j, y in enumerate(b):
                if y:
                    r[i + j] ^= mul(x, y)
    return trim(r)

def p_divmod(a, b):                        # a = q*b + r,  deg r < deg b
    a = trim(a)
    q = [0] * (max(0, len(a) - len(b)) + 1)
    while a != [0] and len(a) >= len(b):
        c = mul(a[-1], inv(b[-1]))
        k = len(a) - len(b)
        q[k] ^= c
        for i, y in enumerate(b):
            a[i + k] ^= mul(c, y)
        a = trim(a)
    return trim(q), a

def p_mod(a, m):
    return p_divmod(a, m)[1]

def p_eval(p, x):
    r = 0
    for c in reversed(p):
        r = mul(r, x) ^ c
    return r

def p_inv_mod(a, m):                       # a^{-1} mod m  (gcd(a,m)=1)
    r0, r1 = trim(m), p_mod(a, m)
    u0, u1 = [0], [1]                      # invariante: r_i = u_i*a mod m
    while len(r1) > 1:                     # hasta llegar a una constante
        assert r1 != [0], "sin inverso: gcd != 1"
        q, r2 = p_divmod(r0, r1)
        r0, r1 = r1, r2
        u0, u1 = u1, p_add(u0, p_mul(q, u1))
    assert r1 != [0], "sin inverso: gcd != 1"
    lam = r1[0]                            # r1 = lam == u1*a  ->  1 = (u1/lam)*a
    return p_mod(p_esc(inv(lam), u1), m)

def p_sqrt_mod(v, g):                      # deg g = 2: anillo GF(256)
    w = p_mod(v, g)                        # raiz = v^(2^7)
    for _ in range(7):
        w = p_mod(p_mul(w, w), g)
    return w

def p_hex(p):
    def t(i):
        return "" if i == 0 else ("x" if i == 1 else f"x^{i}")
    s = '+'.join(f"{c:x}{t(i)}" for i, c in enumerate(p) if c)
    return s or '0'

# ======================================================================
# 3. El codigo de Goppa del juguete (m = 4, n = 16, t = 2)
# ======================================================================

N, M4 = 16, 4
ALPHA = list(range(N))                     # L = (0, 1, z, 1+z, ...)
G_POLY = [0x8, 0x1, 0x1]                   # g(x) = x^2 + x + z^3
INV_G = [inv(p_eval(G_POLY, a)) for a in ALPHA]
COLS = [(INV_G[j], mul(ALPHA[j], INV_G[j])) for j in range(N)]

def sbits(v, n):
    return ''.join(str((v >> i) & 1) for i in range(n))

def vbits(bs):
    return sum(b << i for i, b in enumerate(bs))

def syndrome(word):
    # S(x) = suma_j palabra_j * (x + a_j + 1)/g(a_j)   (mod g)
    S = [0]
    for j, b in enumerate(word):
        if b:
            S = p_add(S, p_mod(p_mul([ALPHA[j] ^ 1, 1], [INV_G[j]]), G_POLY))
    return S

def patterson(S, traza=None):
    R = p_inv_mod(S, G_POLY)                    # R = S^{-1}
    W = p_sqrt_mod(p_add(R, [0, 1]), G_POLY)    # W^2 = x + R
    tau = p_inv_mod(W, G_POLY)                  # tau = W^{-1}
    A, B = p_divmod(G_POLY, tau)                # g = tau*A + B
    sigma = p_add(p_mul(A, A), p_mul([0, 1], p_mul(B, B)))
    if traza is not None:
        traza.update(R=R, W=W, tau=tau, A=A, B=B)
    return sigma

def raices(sigma):
    return [j for j in range(N) if p_eval(sigma, ALPHA[j]) == 0]

# ---- expansion a bits, forma sistematica (I8|T) y G = (T^t|I8) -------

def hhat_variante(variante):
    # variante: (intercambio, invertir_bits)
    swap, msb = variante
    u = lambda j: COLS[j][1] if swap else COLS[j][0]   # filas 0-3
    v = lambda j: COLS[j][0] if swap else COLS[j][1]   # filas 4-7
    H = []
    for r in range(4):
        H.append([(u(j) >> (3 - r if msb else r)) & 1 for j in range(N)])
    for r in range(4):
        H.append([(v(j) >> (3 - r if msb else r)) & 1 for j in range(N)])
    return H

def sistematica(H):
    A = [row[:] for row in H]
    nr = len(A)
    for c in range(nr):
        piv = next((r for r in range(c, nr) if A[r][c]), None)
        if piv is None:
            return None
        A[c], A[piv] = A[piv], A[c]
        for r in range(nr):
            if r != c and A[r][c]:
                A[r] = [x ^ y for x, y in zip(A[r], A[c])]
    return [row[nr:] for row in A]

EXPECT_G = ["1100010110000000", "1100101001000000", "0110110000100000",
            "1001110000010000", "1011101000001000", "0111010100000100",
            "0110101100000010", "1001011100000001"]

def elige_convencion():
    for variante in ((0, 0), (0, 1), (1, 0), (1, 1)):
        H = hhat_variante(variante)
        T = sistematica(H)
        if T is None:
            continue
        G = [sbits(sum(T[a][i] << a for a in range(8)), 8)
             + sbits(1 << i, 8) for i in range(8)]
        if G == EXPECT_G:
            return H, T, G, variante
    raise SystemExit("ninguna convencion reproduce la G del libro")

HHAT, T, G, VARIANTE = elige_convencion()

print("=" * 70)
print("0. GF(16) = F2[z]/(z^4+z+1);  1/g y a/g se calculan en GF(16).")
print("   Potencias de z:", ' '.join(f"z^{e}={EXP[e]:x}" for e in range(15)))

# ----------------------------------------------------------------------
print()
print("=" * 70)
print("1. Las 16 columnas de H~ (tabla del libro):  (1/g(a_j), a_j/g(a_j))")
EXPECT_COLS = {0: (0xf, 0x0), 1: (0xf, 0xf), 2: (0x3, 0x6), 3: (0x3, 0x5),
               4: (0x8, 0x6), 5: (0x8, 0xe), 6: (0x2, 0xc), 7: (0x2, 0xe),
               8: (0xa, 0xf), 9: (0xa, 0x5), 10: (0xc, 0x1), 11: (0xc, 0xd),
               12: (0x5, 0x9), 13: (0x5, 0xc), 14: (0x4, 0xd), 15: (0x4, 0x9)}
for j in range(N):
    a, b = COLS[j]
    assert (a, b) == EXPECT_COLS[j], f"columna {j} discrepa"
    print(f"   j={j:2d}  a_j={ALPHA[j]:x}  g(a_j)={p_eval(G_POLY, ALPHA[j]):x}"
          f"  (1/g, a/g)=({a:x},{b:x})")

# ----------------------------------------------------------------------
print()
print("=" * 70)
print("2. Forma sistematica: filas de G = (T^t | I8)  [T unica]")
for i, row in enumerate(G):
    assert row == EXPECT_G[i]
    print(f"   G_{i} = {row}")
print("   columnas de T (hex):",
      ' '.join(f"{sum(T[a][i] << a for a in range(8)):02x}" for i in range(8)))

# ----------------------------------------------------------------------
print()
print("=" * 70)
print("3. Sindrome del error e = posiciones {10, 13}")
e2 = [0] * N
e2[10] = e2[13] = 1
s0 = 0
s1 = 0
for j in (10, 13):
    s0 ^= COLS[j][0]
    s1 ^= COLS[j][1]
assert (s0, s1) == (0x9, 0xd)              # H~ e = (9, d)
S_e = syndrome(e2)
print(f"   H~ e = (s0, s1) = ({s0:x}, {s1:x})")
print(f"   S(x) = s0*x + (s0+s1) = {p_hex(S_e)}   con S_e = [{S_e[0]:x}, {S_e[1]:x}]")
assert S_e == [0x4, 0x9]

# ----------------------------------------------------------------------
print()
print("=" * 70)
print("4. Patterson paso a paso para S(x) = 4 + 9x (con certificados)")
tr = {}
sig49 = patterson([0x4, 0x9], traza=tr)
R, W, tau, A, B = tr["R"], tr["W"], tr["tau"], tr["A"], tr["B"]
assert R == [0xa, 0x7]
assert p_mul([0x4, 0x9], R) == p_add([1], p_esc(0xa, G_POLY))
assert p_add(R, [0, 1]) == [0xa, 0x6]
assert W == [0xc, 0x7]
assert p_mul(W, W) == p_add([0xa, 0x6], p_esc(0x6, G_POLY))
assert tau == [0x3, 0xb]
assert p_mul(W, tau) == p_add([1], p_esc(0x4, G_POLY))
assert A == [0x3, 0x5] and B == [0xd]
assert p_add(p_mul(tau, A), B) == G_POLY
assert sig49 == [0x5, 0xe, 0x2]
assert raices(sig49) == [10, 13]
q1, r1 = p_divmod(sig49, [0xa, 1])
q2, r2 = p_divmod(sig49, [0xd, 1])
assert r1 == [0] and r2 == [0] and p_add(q1, q2) == [0xe]
assert p_mul(sig49, [0x4, 0x9]) == p_add([0xe], p_mul([0xe, 1], G_POLY))
print(f"   R = S^-1 = {p_hex(R)}        [certificado: S*R = 1 + (z+z^3)*g]")
print(f"   x + R = {p_hex(p_add(R, [0,1]))}")
print(f"   W = raiz(x+R) = {p_hex(W)}   [certificado: W^2 = (x+R) + (z+z^2)*g]")
print(f"   tau = W^-1 = {p_hex(tau)}    [certificado: W*tau = 1 + z^2*g]")
print(f"   g = tau*A + B con A = {p_hex(A)}, B = {p_hex(B)}  (division exacta)")
print(f"   sigma = A^2 + x*B^2 = {p_hex(sig49)}")
print(f"   raices de sigma: {raices(sig49)} = posiciones del error")
print(f"   sigma' = e (0xe) y sigma*S = sigma' + (x+e)*g  [OK]")

# ----------------------------------------------------------------------
print()
print("=" * 70)
print("5. Identidad de van Lint: Gamma2(L,g) = Gamma2(L,g^2)")

def nucleo(H):
    nr = len(H)
    colmask = [sum(H[r][j] << r for r in range(nr)) for j in range(N)]
    tabL, tabH = [0] * 256, [0] * 256
    for b in range(256):
        sL = sH = 0
        for j in range(8):
            if (b >> j) & 1:
                sL ^= colmask[j]
                sH ^= colmask[8 + j]
        tabL[b], tabH[b] = sL, sH
    return frozenset(w for w in range(1 << N)
                     if tabL[w & 255] ^ tabH[w >> 8] == 0)

H4 = []
g2inv = [inv(mul(p_eval(G_POLY, a), p_eval(G_POLY, a))) for a in ALPHA]
for i in range(4):
    for l in range(4):
        H4.append([(mul(1 if i == 0 else pw(ALPHA[j], i), g2inv[j]) >> l) & 1
                   for j in range(N)])
C1, C2 = nucleo(HHAT), nucleo(H4)
assert len(C1) == len(C2) == 256 and C1 == C2
print(f"   |nucleo(H~)| = {len(C1)},  |nucleo(4 filas a^i/g^2)| = {len(C2)}")
print("   las 65536 palabras: coinciden palabra a palabra  [OK]  (y son 256)")

# ----------------------------------------------------------------------
print()
print("=" * 70)
print("6. Correccion: los 120 errores de peso 2 y los 137 sindromes")
ok = 0
for i in range(N):
    for j in range(i + 1, N):
        w = [0] * N
        w[i] = w[j] = 1
        if raices(patterson(syndrome(w))) == [i, j]:
            ok += 1
assert ok == 120
print(f"   Patterson corrige {ok}/120 errores de peso 2  [OK]")
sind = {0}
for j in range(N):
    sind.add((COLS[j][0] << 4) | COLS[j][1])
for i in range(N):
    for j in range(i + 1, N):
        sind.add((COLS[i][0] ^ COLS[j][0]) << 4 | (COLS[i][1] ^ COLS[j][1]))
assert len(sind) == 1 + 16 + 120 == 137
print(f"   de los 2^8 = 256 sindromes, alcanzables con peso <= 2: {len(sind)}"
      "  (1 + 16 + 120)  [OK]")

# ----------------------------------------------------------------------
print()
print("=" * 70)
print("7. McEliece 1978 en miniatura (capitulo 5): G' = S*G*P")

def mmul(A, B):
    return [[sum(A[i][k] * B[k][j] for k in range(len(B))) % 2
             for j in range(len(B[0]))] for i in range(len(A))]

def perm_pull(M, perm):                    # (M P)[i][j] = M[i][perm[j]]
    return [[row[perm[j]] for j in range(len(row))] for row in M]

def inv_f2(M):
    n = len(M)
    A = [row[:] + [1 if i == j else 0 for j in range(n)]
         for i, row in enumerate(M)]
    for c in range(n):
        piv = next(r for r in range(c, n) if A[r][c])
        A[c], A[piv] = A[piv], A[c]
        for r in range(n):
            if r != c and A[r][c]:
                A[r] = [x ^ y for x, y in zip(A[r], A[c])]
    return [row[n:] for row in A]

S8 = [[int(b) for b in row] for row in
      ("01101010", "10010001", "01010001", "01010110",
       "01110110", "00111110", "10011100", "00010101")]
SIG = [14, 2, 7, 10, 4, 8, 12, 6, 9, 15, 11, 5, 0, 3, 13, 1]
SIGINV = [SIG.index(j) for j in range(N)]
Gm = [[int(b) for b in row] for row in EXPECT_G]
GP = mmul(S8, perm_pull(Gm, SIG))
GP_fig = ["1111001110010101", "0000110101111001", "0010000011101001",
          "1000100010100011", "1101000010110010", "1001001000110111",
          "0000011100111110", "0100100101110111"]
for i, row in enumerate(GP):
    txt = sbits(vbits(row), 16)
    assert txt == GP_fig[i], f"G' fila {i} != figura"
print("   G' = S*G*P reproduce la figura del libro  [OK]")

u = [int(c) for c in "01001101"]
z = [int(c) for c in "0001000000001000"]
uGP = [sum(u[i] * GP[i][j] for i in range(8)) % 2 for j in range(N)]
assert sbits(vbits(uGP), 16) == "0000011010001011"
x = [a ^ b for a, b in zip(uGP, z)]
assert sbits(vbits(x), 16) == "0001011010000011"
xPinv = [x[SIGINV[j]] for j in range(N)]
zPinv = [z[SIGINV[j]] for j in range(N)]
assert sbits(vbits(xPinv), 16) == "0100000011101100"
assert sbits(vbits(zPinv), 16) == "1000000000100000"
c = [a ^ b for a, b in zip(xPinv, zPinv)]
assert sbits(vbits(c), 16) == "1100000011001100"
print(f"   u = {sbits(vbits(u), 8)}  (4 unos)   z = {sbits(vbits(z), 16)}")
print(f"   x = u*G' + z = {sbits(vbits(x), 16)}")
print(f"   x*P^-1 = {sbits(vbits(xPinv), 16)}   z*P^-1 = {sbits(vbits(zPinv), 16)}")
print(f"   c = x*P^-1 + z*P^-1 = {sbits(vbits(c), 16)}")
S_x = syndrome(xPinv)
assert S_x == [0x2, 0x3]                   # S(x) = 2 + 3x
sig = patterson(S_x)
assert sig == [0, 0xc, 0xf]                # c x + f x^2
r = raices(sig)
assert r == [0, 10]
print(f"   sindrome: {p_hex(S_x)};  sigma = {p_hex(sig)} (multiplo escalar del localizador)")
print(f"   raices: {r} -> z' recuperado = {sbits(vbits(zPinv), 16)}  [OK]")
uS = c[8:]
u_rec = [sum(uS[i] * inv_f2(S8)[i][k] for i in range(8)) % 2 for k in range(8)]
assert sbits(vbits(uS), 8) == "11001100"
assert u_rec == u
print(f"   u*S = {sbits(vbits(uS), 8)};  u = (u*S)*S^-1 = {sbits(vbits(u_rec), 8)}"
      "  -> mensaje recuperado  [OK]")

# ----------------------------------------------------------------------
print()
print("=" * 70)
print("8. KEM del juguete (capitulo 11)")

DELTA = bytes([0x5a]) * 32
flujo = shake_256(bytes([0x40]) + DELTA).digest(8)   # etiqueta 64 = 0x40
s_val = flujo[0] | (flujo[1] << 8)
assert s_val == 0xCBAE
s_str = sbits(s_val, 16)
print(f"   delta = 5a repetido 32 veces")
print(f"   s = {s_str[:4]} {s_str[4:8]} {s_str[8:12]} {s_str[12:]}"
      f" = {s_val:04x}  (bytes AE CB, little-endian)")

e_hi, e_lo = 0x24, 0x00                    # e = 0x2400: posiciones 10 y 13
C_bits = [(e_lo >> r) & 1 ^
          sum(T[r][j] * ((e_hi >> j) & 1) for j in range(8)) % 2
          for r in range(8)]
C_val = vbits(C_bits)
assert C_val == 0x98
print(f"   e = 0000000000100100 (posiciones 10 y 13)")
print(f"   C = e_alto + T*(e_bajo) = col2(T) + col5(T) = 36 ^ AE = {C_val:02x}"
      f" = {sbits(C_val, 8)}")

K = shake_256(bytes([0x01]) + bytes([e_lo, e_hi]) + bytes([C_val])).digest(32)
assert K.hex() == \
    "597f46fbba325c68f0b0b45c6ea984961461eb220d1f24024ed888912b336278"
print(f"   K  = SHAKE256(01|00 24|98) = {K.hex()}")

v = [(C_val >> r) & 1 for r in range(8)] + [0] * 8
S_v = syndrome(v)
assert S_v == [0x4, 0x9]
sig_v = patterson(S_v)
assert sig_v == [0x5, 0xe, 0x2]
r_v = raices(sig_v)
assert r_v == [10, 13]
e_dec = [0] * N
for p in r_v:
    e_dec[p] = 1
assert e_dec == e2
K2 = shake_256(bytes([0x01]) + bytes([e_lo, e_hi]) + bytes([C_val])).digest(32)
assert K2 == K
print(f"   decap: v = (C, 0); S(v) = {p_hex(S_v)}; sigma = {p_hex(sig_v)}")
print(f"   raices {r_v} -> e correcto;  K' = K  [OK]")

C_t = 0x99
v_t = [(C_t >> r) & 1 for r in range(8)] + [0] * 8
S_t = syndrome(v_t)
sig_t = patterson(S_t)
assert raices(sig_t) == []
K2t = shake_256(bytes([0x00]) + bytes([flujo[0], flujo[1]]) + bytes([C_t])).digest(32)
assert K2t.hex() == \
    "ba1a72ac26b0fcf20c55c108bb681e27aa9c270c60802a821601130b7f1e1494"
print(f"   C~ = C+1 = {sbits(C_t, 8)}: S~ = {p_hex(S_t)}, sigma~ = {p_hex(sig_t)},"
      " sin raices")
print(f"   K'' = SHAKE256(00|AE CB|99) = {K2t.hex()}  (rechazo)")

print()
print("=" * 70)
print("TODAS LAS COMPROBACIONES OK")
\end{lstlisting}

\section{La salida}

La salida siguiente es la de una ejecución real (octubre de 2026),
sin más recorte que los espacios en blanco finales de línea:

\begin{lstlisting}[numbers=none]
======================================================================
0. GF(16) = F2[z]/(z^4+z+1);  1/g y a/g se calculan en GF(16).
   Potencias de z: z^0=1 z^1=2 z^2=4 z^3=8 z^4=3 z^5=6 z^6=c z^7=b z^8=5 z^9=a z^10=7 z^11=e z^12=f z^13=d z^14=9

======================================================================
1. Las 16 columnas de H~ (tabla del libro):  (1/g(a_j), a_j/g(a_j))
   j= 0  a_j=0  g(a_j)=8  (1/g, a/g)=(f,0)
   j= 1  a_j=1  g(a_j)=8  (1/g, a/g)=(f,f)
   j= 2  a_j=2  g(a_j)=e  (1/g, a/g)=(3,6)
   j= 3  a_j=3  g(a_j)=e  (1/g, a/g)=(3,5)
   j= 4  a_j=4  g(a_j)=f  (1/g, a/g)=(8,6)
   j= 5  a_j=5  g(a_j)=f  (1/g, a/g)=(8,e)
   j= 6  a_j=6  g(a_j)=9  (1/g, a/g)=(2,c)
   j= 7  a_j=7  g(a_j)=9  (1/g, a/g)=(2,e)
   j= 8  a_j=8  g(a_j)=c  (1/g, a/g)=(a,f)
   j= 9  a_j=9  g(a_j)=c  (1/g, a/g)=(a,5)
   j=10  a_j=a  g(a_j)=a  (1/g, a/g)=(c,1)
   j=11  a_j=b  g(a_j)=a  (1/g, a/g)=(c,d)
   j=12  a_j=c  g(a_j)=b  (1/g, a/g)=(5,9)
   j=13  a_j=d  g(a_j)=b  (1/g, a/g)=(5,c)
   j=14  a_j=e  g(a_j)=d  (1/g, a/g)=(4,d)
   j=15  a_j=f  g(a_j)=d  (1/g, a/g)=(4,9)

======================================================================
2. Forma sistematica: filas de G = (T^t | I8)  [T unica]
   G_0 = 1100010110000000
   G_1 = 1100101001000000
   G_2 = 0110110000100000
   G_3 = 1001110000010000
   G_4 = 1011101000001000
   G_5 = 0111010100000100
   G_6 = 0110101100000010
   G_7 = 1001011100000001
   columnas de T (hex): a3 53 36 39 5d ae d6 e9

======================================================================
3. Sindrome del error e = posiciones {10, 13}
   H~ e = (s0, s1) = (9, d)
   S(x) = s0*x + (s0+s1) = 4+9x   con S_e = [4, 9]

======================================================================
4. Patterson paso a paso para S(x) = 4 + 9x (con certificados)
   R = S^-1 = a+7x        [certificado: S*R = 1 + (z+z^3)*g]
   x + R = a+6x
   W = raiz(x+R) = c+7x   [certificado: W^2 = (x+R) + (z+z^2)*g]
   tau = W^-1 = 3+bx    [certificado: W*tau = 1 + z^2*g]
   g = tau*A + B con A = 3+5x, B = d  (division exacta)
   sigma = A^2 + x*B^2 = 5+ex+2x^2
   raices de sigma: [10, 13] = posiciones del error
   sigma' = e (0xe) y sigma*S = sigma' + (x+e)*g  [OK]

======================================================================
5. Identidad de van Lint: Gamma2(L,g) = Gamma2(L,g^2)
   |nucleo(H~)| = 256,  |nucleo(4 filas a^i/g^2)| = 256
   las 65536 palabras: coinciden palabra a palabra  [OK]  (y son 256)

======================================================================
6. Correccion: los 120 errores de peso 2 y los 137 sindromes
   Patterson corrige 120/120 errores de peso 2  [OK]
   de los 2^8 = 256 sindromes, alcanzables con peso <= 2: 137  (1 + 16 + 120)  [OK]

======================================================================
7. McEliece 1978 en miniatura (capitulo 5): G' = S*G*P
   G' = S*G*P reproduce la figura del libro  [OK]
   u = 01001101  (4 unos)   z = 0001000000001000
   x = u*G' + z = 0001011010000011
   x*P^-1 = 0100000011101100   z*P^-1 = 1000000000100000
   c = x*P^-1 + z*P^-1 = 1100000011001100
   sindrome: 2+3x;  sigma = cx+fx^2 (multiplo escalar del localizador)
   raices: [0, 10] -> z' recuperado = 1000000000100000  [OK]
   u*S = 11001100;  u = (u*S)*S^-1 = 01001101  -> mensaje recuperado  [OK]

======================================================================
8. KEM del juguete (capitulo 11)
   delta = 5a repetido 32 veces
   s = 0111 0101 1101 0011 = cbae  (bytes AE CB, little-endian)
   e = 0000000000100100 (posiciones 10 y 13)
   C = e_alto + T*(e_bajo) = col2(T) + col5(T) = 36 ^ AE = 98 = 00011001
   K  = SHAKE256(01|00 24|98) = 597f46fbba325c68f0b0b45c6ea984961461eb220d1f24024ed888912b336278
   decap: v = (C, 0); S(v) = 4+9x; sigma = 5+ex+2x^2
   raices [10, 13] -> e correcto;  K' = K  [OK]
   C~ = C+1 = 10011001: S~ = b+6x, sigma~ = 5+8x+ax^2, sin raices
   K'' = SHAKE256(00|AE CB|99) = ba1a72ac26b0fcf20c55c108bb681e27aa9c270c60802a821601130b7f1e1494  (rechazo)

======================================================================
TODAS LAS COMPROBACIONES OK
\end{lstlisting}

\section{Simplificaciones del juguete}

El juguete no implementa la especificación real: la ilustra. Sus
diferencias con Classic McEliece, todas deliberadas:

\begin{itemize}[leftmargin=1.5em,itemsep=3pt]
  \item \textbf{El error es fijo.} \textsf{FixedWeight} devuelve
        siempre el error de soporte $\{10,13\}$: alcanza para
        verificar la tubería, pero la especificación lo sortea
        uniformemente entre las $\binom{n}{t}$ posibilidades
        (proposición~\ref{prop:fixedweight-uniform}).
  \item \textbf{No hay semillas, ni reintentos, ni variantes.} El
        soporte, $g$ y $\sigma$ están fijados a mano: no aparecen
        $\sigma_1$ y $\sigma_2$, ni la forma semi-sistemática
        $(\mu,\nu)=(32,64)$, ni la red de Benes que comprime
        $\alpha$ en la clave secreta (todo el
        capítulo~\ref{ch:keygen}); tampoco las variantes
        \texttt{f}, \texttt{pc} y \texttt{pcf}
        (capítulo~\ref{ch:parametros}).
  \item \textbf{La seguridad es nula.} Con $n=16$, enumerar las
        $2^{16}$ palabras es un parpadeo: el juguete verifica la
        corrección, nunca la dificultad. El problema «duro» empieza
        en los capítulos~\ref{ch:isd} y~\ref{ch:estructura}.
  \item \textbf{Las etiquetas del hash sí son las de la
        especificación}: \texttt{0x40} para el generador
        $G(\delta)$, \texttt{0x01} y \texttt{0x00} para las dos
        ramas de la clave de sesión (capítulo~\ref{ch:labkem}), con
        la convención de bits LSB-primero y bytes little-endian de
        la sección~6.2 de~\cite{cmspec}.
\end{itemize}

\medskip
\noindent Para las cuentas a mano, el apéndice~\ref{app:gf16} tiene
las tablas de $\FF4$ en formato de bolsillo.
